\section{Functors}\label{sec:functors}

\paragraph{Functors}

\begin{definition}\label{def:functor}\mcite[13]{MacLane1998CategoryTheory}
  A \term[ru=функтор (\cite[def. III.1.1]{ЦаленкоШульгейфер1974ОсновыТеорииКатегорий})]{functor} \( F: \cat{C} \to \cat{D} \) between the \hyperref[def:category]{categories} \( \cat{C} \) and \( \cat{D} \) is a pair of functions
  \begin{equation*}
    \begin{aligned}
      F_{\obj}: &\obj(\cat{C}) \to \obj(\cat{D}), \\
      F_{\mor}: &\mor(\cat{C}) \to \mor(\cat{D}),
    \end{aligned}
  \end{equation*}
  such that
  \begin{equation*}
    F_{\mor}(f: A \to B): F_{\obj}(A) \to F_{\obj}(B)
  \end{equation*}
  and with compatibility conditions for identities and composition.

  We will use the same notation for both functions, i.e. we write \( F(A) \) for an object \( A \) and \( F(f) \) for a morphism \( f \). This is imprecise because, as per \cref{rem:disjoint_graph_vertices_and_arcs}, we have not assumed that \( F_{\obj} \) and \( F_{\mor} \) are disjoint, and thus \( F(A) \) and \( F(f) \) would differ even if \( A = f \). Nevertheless, this notation reflects popular convention, so we stick to it.

  It remains to state the compatibility conditions, which are collectively called \term[en=functoriality axioms (\cite[def. 1.3.1]{Perrone2024StartingCategoryTheory})]{functoriality}:
  \begin{thmenum}[resume=def:functor]
    \thmitem[def:functor/CF1]{CF1} Functors must preserve identities: for any object \( A \) in \( \cat{C} \), the following equality must hold:
    \begin{equation*}
      F(\id_A) = \id_{F(A)}.
    \end{equation*}

    \thmitem[def:functor/CF2]{CF2} Functors must preserve composition: for any pair of morphism \( f: A \to B \) and \( g: B \to C \) in \( \cat{C} \), the following equality must hold:
    \begin{equation*}
      F(g \bincirc f) = F(g) \bincirc F(f).
    \end{equation*}
  \end{thmenum}
\end{definition}
\begin{comments}
  \item As mentioned in \cref{rem:category_as_graph}, the notion of functor precedes that of category, and is attributed to \cite{EilenbergMacLane1942NaturalIsomorphisms}. The paper establishes \ref{def:functor/CF1} and \ref{def:functor/CF2}, as well as using the same symbol for both the object and arrow mapping.

  \item There is a closely related notion called a \enquote{contravariant functor} --- see \cref{def:contravariant_functor}.
\end{comments}

\begin{proposition}\label{thm:functor_is_graph_homomorphism}
  If, as per \cref{def:category}, we regard categories as \hyperref[def:directed_multigraph]{directed multigraphs}, the \hyperref[def:functor]{functors} are \hyperref[def:directed_multigraph/morphism]{graph homomorphism} (but not vice versa).
\end{proposition}
\begin{proof}
  Given a functor \( F: \cat{C} \to \cat{D} \), for any morphism \( f: A \to B \) in \( \cat{C} \), we have
  \begin{equation*}
    F(\dom(f)) = F(A) = \dom(F(f)),
  \end{equation*}
  which implies \eqref{eq:def:directed_multigraph/morphism/head}. We also have
  \begin{equation*}
    F(\co\dom(f)) = F(B) = \co\dom(F(f)),
  \end{equation*}
  which implies \eqref{eq:def:directed_multigraph/morphism/tail}.
\end{proof}

\begin{definition}\label{def:endofunctor}\mimprovised
  We call a \hyperref[def:functor]{functor} from a category to itself an \term{endomorphism}.
\end{definition}

\begin{example}\label{ex:def:functor}
  We list examples of \hyperref[def:functor]{functors}.

  \begin{thmenum}
    \thmitem{ex:def:functor/power_set} The \hyperref[def:basic_set_operations/power_set]{power set} \( \pow: \cat{Set} \to \cat{Set} \) is a canonical example of an \hyperref[def:endofunctor]{endofunctor}. Explicitly:
    \begin{equation*}
      \begin{aligned}
        &\pow: \cat{Set} \to \cat{Set}, \\
        &\pow(A) \coloneqq \set{ S \given S \subseteq A }, \\
        &\pow(f: A \to B) \coloneqq (S \mapsto f[S]). \\
      \end{aligned}
    \end{equation*}

    The nuance here is that the functor sends every function \( f: A \to B \) to another function, which takes a subset \( S \) of \( A \) to its \hyperref[def:set_valued_map/image]{set value} \( f[S] \).

    We must verify that it is indeed a functor. \ref{def:functor/CF2} is satisfied because
    \begin{equation*}
      \pow(g) \bincirc \pow(f) = (S \mapsto g[f[S]]) = \pow(g \bincirc f).
    \end{equation*}

    The condition \ref{def:functor/CF1} is also obviously satisfied.

    \thmitem{ex:def:functor/image_no_category}\mcite{MathSE:image_of_functor_is_not_a_category} If, as per \cref{thm:functor_is_graph_homomorphism}, we regard a functor as a \hyperref[def:directed_multigraph/morphism]{graph homomorphism}, we obtain some notion of an image that generalizes \hyperref[def:set_valued_map/image]{images of functions}.

    This image is a directed multigraph, but in general not a category. As an example, consider the functor \( F: \cat{C} \to \cat{D} \) from \cref{fig:ex:def:functor/image_no_category}.

    \begin{figure}
      \hfill
      \includegraphics[page=1]{output/ex__def__functor}
      \hfill
      \hfill
      \caption{A functor whose image is not a category.}\label{fig:ex:def:functor/image_no_category}
    \end{figure}

    \begin{itemize}
      \item The solid arrows are the morphisms in \( \cat{C} \) and their images in \( F(\cat{C}) \).
      \item The dashed arrows denote the action of the functor \( F \).
      \item The dotted arrow exists in \( \cat{D} \) as the composition of the other two arrows, however it is missing in the image \( F(\cat{C}) \). Thus, composition is not fully defined in \( F(\cat{C}) \), and \( F(\cat{C}) \) fails to be a category.
    \end{itemize}
  \end{thmenum}
\end{example}

\begin{proposition}\label{thm:def:functor}
  \hyperref[def:functor]{Functors} have the following basic properties:
  \begin{thmenum}
    \thmitem{thm:def:functor/one_sided_inverses} Functors preserve one-sided inverses: for every functor \( F: \cat{C} \to \cat{D} \) and every morphism \( f: A \to B \) in \( \cat{C} \) with a left inverse \( g: B \to A \), \( F(f) \) is a left inverse of \( F(g) \).

    Equivalently, if \( g \) is a right inverse of \( f \), then \( F(g) \) is a right inverse of \( F(f) \).

    \thmitem{thm:def:functor/inverses} Functors preserve two-sided inverses: For every functor \( F: \cat{C} \to \cat{D} \) and every isomorphism \( f: A \to B \) in \( \cat{C} \), \( F(f^{-1}) \) is the inverse of \( F(f) \).

    In particular, if \( f \) is an \hyperref[def:morphism_invertibility/isomorphism]{isomorphisms}, so is \( F(f) \), but the converse only holds for some functors --- see \cref{thm:def:functor_invertibility/fully_faithful_reflects_invertible}.
  \end{thmenum}
\end{proposition}
\begin{proof}
  \SubProofOf{thm:def:functor/one_sided_inverses} Let \( g: B \to A \) be a left inverse of \( f: A \to B \) in \( \cat{C} \). Then
  \begin{equation*}
    F(g) \bincirc F(f)
    \reloset {\ref{def:functor/CF2}} =
    F(g \bincirc f)
    =
    F(\id_A)
    \reloset {\ref{def:functor/CF1}} =
    \id_{F(A)}.
  \end{equation*}

  Thus, \( F(f) \) is a right inverse of \( F(g) \). Since \( g \) is a left inverse of \( f \), automatically \( F(g) \) is a left inverse of \( F(f) \).

  \SubProofOf{thm:def:functor/inverses} If \( f^{-1} \) is a two-sided inverse of \( f \), \cref{thm:def:functor/one_sided_inverses} implies that \( F(f^{-1}) \) is both a left and right inverse of \( F(f) \).
\end{proof}

\begin{definition}\label{def:contravariant_functor}\mcite[33]{MacLane1998CategoryTheory}
  In \cref{def:functor}, we have required a functor \( F: \cat{C} \to \cat{D} \) to obey \ref{def:functor/CF2}:
  \begin{equation*}
    F(g \bincirc f) = F(g) \bincirc F(f).
  \end{equation*}

  \begin{thmenum}
    \thmitem[def:contravariant_functor/CF2]{CF2\( ^\oppos \)} If we take the \hyperref[def:opposite_category]{opposite category} \( \cat{C}^\oppos \) instead of \( \cat{C} \) (or \( \cat{D}^\oppos \) instead of \( \cat{D} \)), this condition becomes
    \begin{equation*}
      F(g \bincirc f) = F(f) \bincirc F(g).
    \end{equation*}
  \end{thmenum}

  We call this a \term[ru=контравариантный функтор (\cite[def. III.1.1]{ЦаленкоШульгейфер1974ОсновыТеорииКатегорий})]{contravariant functor} from \( \cat{C} \) to \( \cat{D} \), and we call the usual notion of functor \term[ru=ковариантный (функтор) (\cite[def. III.1.1]{ЦаленкоШульгейфер1974ОсновыТеорииКатегорий})]{covariant}.
\end{definition}
\begin{comments}
  \item Since a contravariant functor is by definition a covariant functor, to avoid confusion, we only use the term \enquote{contravariant} sparingly.

  \item This definition is an instance of the \hyperref[rem:red_herring_principle]{mathematical red herring principle}.
\end{comments}

\begin{example}\label{ex:def:contravariant_functor}
  We list examples of \hyperref[def:contravariant_functor]{contravariant functors}:
  \begin{thmenum}
    \thmitem{ex:def:contravariant_functor/dual_space}\mcite[33]{MacLane1998CategoryTheory} We can try to na\"ively define a functor that assigns to a \hyperref[def:vector_space]{vector space} its \hyperref[def:dual_vector_space]{algebraic dual}:
    \begin{equation*}
      \begin{aligned}
        &F: \cat[\BbbK]{Vect} \to \cat[\BbbK]{Vect}, \\
        &F(V) \coloneqq V^*, \\
        &F(f: V \to W) \coloneqq (\varphi: W \to \BbbK \mapsto \varphi \bincirc f).
      \end{aligned}
    \end{equation*}

    However, \( F(f) \) is actually a function from \( W^* \) to \( V^* \). This makes \( F \) a contravariant functor or, equivalently, a functor from \( \cat[\BbbK]{Vect}^\oppos \) to \( \cat[\BbbK]{Vect} \).
  \end{thmenum}
\end{example}

\paragraph{Category of small categories}

\begin{proposition}\label{thm:small_functors}
  A \hyperref[def:functor]{functor} \( F: \cat{C} \to \cat{D} \), regarded as a pair \( (F_{\obj}, F_{\mor}) \) of functions, is a \hyperref[def:small_set]{small set} if and only if \( F_{\mor} \) is small (under both interpretation of of function smallness from \cref{rem:small_functions}).
\end{proposition}
\begin{comments}
  \item The general ambiguity of the term \enquote{small functor} is discussed in \cref{rem:small_functors}.

  Under both interpretations, if \( \cat{C} \) and \( \cat{D} \) are small categories, then \( F \) is a small set.
\end{comments}
\begin{proof}
  \SufficiencySubProof* If \( F = (F_{\obj}, F_{\mor}) \) is a small set, \cref{thm:small_set_closure/indexed_family} implies that so is \( F_{\mor} \).

  \NecessitySubProof* Conversely, suppose that \( F_{\mor} \) is a small set.

  \begin{itemize}
    \item Under the interpretation from \cref{rem:small_functions/triple}, where \( F_{\mor} \) is the triple
    \begin{equation*}
      (\gr F_{\mor}, \mor(\cat{C}), \mor(\cat{D})),
    \end{equation*}
    since \( F_{\mor} \) is assumed to be small, so are \( \mor(\cat{C}) \) and \( \mor(\cat{D}) \) are. \Cref{thm:small_set_closure/category} implies that \( \obj(\cat{C}) \) and \( \obj(\cat{D}) \) are also small.

    Thus, the domain and codomain of \( F_{\obj} \) are small, and hence \( F_{\obj} \) itself is small.

    Then \( F = (F_{\obj}, F_{\mor}) \) is also a small set by \cref{thm:small_set_closure/indexed_family}.

    \item Under the interpretation from \cref{rem:small_functions/graph}, where \( F_{\mor} \) is its graph \( \gr F_{\mor} \), since \( F_{\mor} \) is assumed to be small, so are \( \mor(\cat{C}) \) and the image \( \img F_{\mor} \).

    By the same reasoning as above, we conclude that \( F_{\obj} \) is also small, and so is \( F = (F_{\obj}, F_{\mor}) \).
  \end{itemize}
\end{proof}

\begin{remark}\label{rem:small_functors}
  Depending on how we encode the \hyperref[def:functor]{functor} \( F: \cat{C} \to \cat{D} \) as a set, there are incompatible ways in which it can be a \hyperref[def:small_set]{small set}, i.e. a member of the \hyperref[def:smallest_grothendieck_universe]{smallest Grothendieck universe} \( \mscrU_0 \). This is an extension of our discussion of small functions from \cref{rem:small_functions}.

  Due to the ambiguity, we will generally avoid speaking about \enquote{small functors}. We do, however, need to consider small \hyperref[def:categorical_diagram]{diagrams}, so we adjust their definition in \cref{def:small_diagram} specifically so that they can be true small sets. There are many examples where the size of a diagram is essential. For example, a \hyperref[def:complete_category]{complete category} is defined as one in which all limits of small diagrams exist.

  \begin{thmenum}
    \thmitem{rem:small_functors/triple} One options is to regard \( F: \cat{C} \to \cat{D} \) as a pair \( F = (F_{\obj}, F_{\mor}) \) of functions, and, as per \cref{rem:function_definition}, then regard these functions as the triples
    \begin{align*}
      \parens[\big]{ \gr F_{\obj}, \obj(\cat{C}), \obj(\cat{D}) }, && \parens[\big]{ \gr F_{\mor}, \mor(\cat{C}), \mor(\cat{D}) }.
    \end{align*}

    Then \cref{thm:small_set_closure/function} implies that \( F \) is a small set if and only if both \( \cat{C} \) and \( \cat{D} \) are \hyperref[def:small_category]{small categories}.

    Unfortunately, this makes it impossible to define a small functor into a large category like \( \cat{Set} \). This is the reason we do not treat diagrams as functors in \cref{def:categorical_diagram} --- we want to have a notion of \enquote{small diagram} even though the more general \enquote{small functor} is generally ambiguous.

    \thmitem{rem:small_functors/graph} Alternatively, we can define \( F \) as the pair \( (\gr F_{\obj}, \gr F_{\mor}) \) so that \( F \) is a small set if and only if \( \cat{C} \) and its image in \( \cat{D} \) are small categories.

    This makes it possible for a functor into \( \cat{Set} \) to itself be small. Since a diagram can be regarded as a functor, small diagrams simply becomes small functors.

    As noted in \cref{rem:small_category/reflection}, \textcite[26]{Shulman2008SetTheoryForCategoryTheory} discusses how in the axiomatization \logic{ZFC/S}, where a weaker notion of universe is used instead of Grothendieck universes, a functor from a small category to \( \cat{Set} \) may fail to be a small set, and so in some cases it makes sense to have a dedicated notion of small functors.
  \end{thmenum}
\end{remark}

\begin{definition}\label{def:category_of_small_categories}\mcite[14]{MacLane1998CategoryTheory}
  We define the \term{category of small categories} \( \cat{Cat} \) as follows:
  \begin{thmenum}
    \thmitem{def:category_of_small_categories/objects} The \hyperref[def:category/objects]{objects} are, unsurprisingly, the \hyperref[def:small_category]{small category}, i.e. the set of all categories \( \cat{C} \) whose set of morphisms \( \mor(\cat{C}) \) belongs to the \hyperref[def:smallest_grothendieck_universe]{smallest Grothendieck universe} \( \mscrU_0 \).

    \thmitem{def:category_of_small_categories/morphisms} The \hyperref[def:category/morphisms]{morphisms} between small categories are the \hyperref[def:functor]{functors} between them.

    Note that, by \cref{thm:small_functors}, functors between small category are themselves small as sets (under both interpretations of functor smallness from \cref{rem:small_functors}).

    \thmitem{def:category_of_small_categories/composition} The \hyperref[def:category/composition]{composition operation} is defined for \( F: \cat{C} \to \cat{D} \) and \( G: \cat{D} \to \cat{E} \) componentwise, same as the composition of directed multigraph homomorphisms from \cref{def:directed_multigraph/category/composition}.

    \thmitem{def:category_of_small_categories/identity} The \hyperref[def:category/identity]{identity functor} on \( \cat{C} \) is simply the pair of \hyperref[def:set_valued_map/identity]{identity function}
    \begin{equation*}
      \id_{\cat{C}} = \parens[\big]{ \id_{\obj(\cat{C})}, \id_{\mor(\cat{C})} }.
    \end{equation*}
  \end{thmenum}
\end{definition}
\begin{defproof}
  To see that \( \cat{Cat} \) is indeed a category, we verify the conditions \ref{def:category/C1} and \ref{def:category/C2}.

  \SubProofOf{def:category/C1} For every pair of small categories \( \cat{C} \) and \( \cat{D} \) and every functor \( F: \cat{C} \to \cat{D} \), for every object \( A \) in \( \cat{C} \) we have
  \begin{equation*}
    [\id_{\cat{D}} \bincirc F](A)
    =
    \id_{\cat{D}}(F(A))
    =
    F(A)
    =
    F(\id_{\cat{C}}(A))
    =
    [F \bincirc \id_{\cat{C}}](A)
  \end{equation*}
  and analogously for morphisms.

  Therefore, \( \id_{\cat{C}} \) and \( \id_{\cat{D}} \) satisfy \ref{def:category/C1}.

  \SubProofOf{def:category/C2} Associativity of functor composition is inherited from the associativity of function composition.
\end{defproof}

\begin{definition}\label{def:subcategory}\mcite[15]{MacLane1998CategoryTheory}
  A \term[ru=подкатегория (\cite[15]{ЦаленкоШульгейфер1974ОсновыТеорииКатегорий})]{subcategory} of \( \cat{C} \) is a \hyperref[def:directed_multigraph/subobject]{subgraph} of \( \cat{C} \) that is closed under identities and composition.

  \begin{thmenum}
    \thmitem{def:subcategory/inclusion} For every subcategory \( \cat{D} \) of \( \cat{C} \) there exists an \term{inclusion functor} \( \Iota: \cat{D} \to \cat{C} \), which sends every object and morphism of \( \cat{D} \) to itself in \( \cat{C} \).

    \thmitem{def:subcategory/full} We say that the subcategory \( \cat{D} \) is \term[ru=полная (подкатегория) (\cite[15]{ЦаленкоШульгейфер1974ОсновыТеорииКатегорий})]{full} if it contains all morphisms from \( \cat{C} \) whose domain and codomain is in \( \cat{D} \), i.e. if it is an \hyperref[def:induced_subgraph]{induced subgraph}.

    By \cref{thm:def:functor_invertibility/full_subcategory}, this is equivalent to the inclusion functor being \hyperref[def:functor_invertibility/full]{full}.
  \end{thmenum}
\end{definition}

\begin{proposition}\label{thm:def:subcategory}
  \hyperref[def:subcategory]{Subcategories} have the following basic properties:
  \begin{thmenum}
    \thmitem{thm:def:subcategory/subobject} The subcategories of a small category \( \cat{C} \) correspond exactly to the \hyperref[def:categorical_subobject]{subobjects} of \( \cat{C} \) in \( \cat{Cat} \). More specifically, if \( \cat{D} \) and \( \cat{E} \) are subcategories of \( \cat{C} \), the inclusions \( \Iota_{\cat{D}}: \cat{D} \to \cat{C} \) and \( \iota_{\cat{E}}: \cat{E} \to \cat{C} \) induce the same \hyperref[def:categorical_subobject]{categorical subobjects} of \( \cat{C} \) if and only if \( \cat{D} = \cat{E} \).
  \end{thmenum}
\end{proposition}
\begin{proof}
  \SubProofOf{thm:def:subcategory/subobject} Follows from the analogous result for \hyperref[def:directed_multigraph]{directed multigraphs} shown in \cref{thm:subgraphs}.
\end{proof}

\begin{definition}\label{def:empty_category}\mcite[10]{MacLane1998CategoryTheory}
  There is an obvious \hyperref[def:category]{category} with no objects. We call it the \term[ru=пустая категория (\cite[16]{ЦаленкоШульгейфер1974ОсновыТеорииКатегорий})]{empty category} and denote it by \( \cat{0} \).
\end{definition}
\begin{comments}
  \item This is an \hyperref[def:universal_objects/initial]{initial object} in the \hyperref[def:category_of_small_categories]{category of small categories} --- see \cref{thm:universal_categories}.
\end{comments}

\begin{definition}\label{def:discrete_category}\mcite[11]{MacLane1998CategoryTheory}
  A \term[ru=дискретная (категория) (\cite[15]{ЦаленкоШульгейфер1974ОсновыТеорииКатегорий})]{discrete category} is a category with no morphisms except for the identities.
\end{definition}
\begin{comments}
  \item Clearly to any set there corresponds exactly one discrete category and vice versa.

  \item The concept is based on \hyperref[def:discrete_topology]{discrete topologies} and is unrelated to Aristotelian discrete variables in the sense of \cref{con:discrete_and_continuous_variables}.
\end{comments}

\begin{example}\label{ex:discrete_category_adjunction}
  Denote by
  \begin{equation*}
    U: \cat{Cat} \to \cat{Set}
  \end{equation*}
  the functor that for any small category \( \cat{C} \) gives us its set of objects \( \obj(\cat{C}) \). There is also a functor
  \begin{equation*}
    D: \cat{Set} \to \cat{Cat}
  \end{equation*}
  that for any small set \( A \) gives us the \hyperref[def:discrete_category]{discrete category} whose set of objects is \( A \).

  This is actually an \hyperref[def:category_adjunction]{adjunction} --- see \cref{ex:def:category_adjunction/set_cat}.
\end{example}

\begin{definition}\label{def:singleton_category}\mimprovised
  We refer to the \hyperref[def:discrete_category]{discrete category} of a \hyperref[def:singleton_set]{singleton set} \( \set{ A } \) as a \term{singleton category}.

  As shown in \cref{thm:universal_categories}, the singleton category of a \hyperref[def:small_set]{small set} is a \hyperref[def:universal_objects/terminal]{terminal object} in the \hyperref[def:category_of_small_categories]{category of small categories}, so it is also called a \term[en=terminal category (\cite[275]{MacLane1998CategoryTheory})]{terminal category}.

  By \cref{thm:def:universal_objects/unique}, a terminal category is unique up to a unique isomorphism. If we need a fixed representative, we can use the singleton category of \( \varnothing \), which we will denote by \( \cat{1} \).
\end{definition}

\begin{proposition}\label{thm:universal_categories}
  In the \hyperref[def:category_of_small_categories]{category of small categories}, the \hyperref[def:empty_category]{empty category} is an \hyperref[def:universal_objects/initial]{initial object}, while the \hyperref[def:singleton_category]{singleton category} of any \hyperref[def:small_set]{small set} is a \hyperref[def:universal_objects/terminal]{terminal object}.
\end{proposition}
\begin{proof}
  Fix a small category \( \cat{C} \). The empty functor is the unique morphism from the empty category \( 0 \) to \( \cat{C} \), while the unique morphism from \( \cat{C} \) to a singleton category \( \cat{S} \) sends every object and morphism in \( \cat{C} \) to the the unique object and morphism in \( \cat{S} \).
\end{proof}

\paragraph{Categorical diagrams}

\begin{definition}\label{def:categorical_diagram}
  Fix a category \( \cat{I} \), called an \term{index category}. A \term{diagram} in \( \cat{C} \) of shape \( \cat{I} \) is simply a functor \( D: \cat{I} \to \cat{C} \), whose domain is \( \cat{I} \). We sometimes identify a diagram functor with its image \( D(\cat{I}) \).

  It is often convenient to draw graphically the \hyperref[def:graph_geometric_realization]{geometric realizations} of the \hyperref[def:directed_multigraph]{multigraph} \( D(\cat{I}) \). An established convention is to allow multiple vertices representing the same object, which can be achieved formally by actually adjoining new vertices to the graph and labeling them as in \cref{def:labeled_set}. Other established conventions for drawing diagrams include not drawing identity morphisms and adding various visual aids. In this regard, categorical diagrams correspond to the everyday sense of the word \enquote{diagram}.

  We say that the diagram is \term{small} if its index category is \hyperref[def:small_category]{small}.

  We say that the diagram \( D \) over \( \cat{C} \) \term{commutes} if, whenever
  \begin{equation*}
    A \reloset {f_1} \to \anon \reloset {f_2} \to \cdots \reloset {f_{n-1}} \to \anon \reloset {f_n} \to B
  \end{equation*}
  and
  \begin{equation*}
    A \reloset {g_1} \to \anon \reloset {g_2} \to \cdots \reloset {g_{n-1}} \to \anon \reloset {g_m} \to B
  \end{equation*}
  are two \hyperref[def:graph_walk/directed]{directed walks} in the graph \( D(\cat{I}) \) with identical endpoints and either \( n > 1 \) or \( m > 1 \), then
  \begin{equation*}
    f_n \bincirc f_{n-1} \bincirc \cdots \bincirc f_2 \bincirc f_1
    =
    g_m \bincirc g_{m-1} \bincirc \cdots \bincirc g_2 \bincirc g_1.
  \end{equation*}

  We do not really care about how the objects and morphisms in \( \cat{I} \) are labeled, hence we often use placeholder dots like in \eqref{eq:ex:directed_multigraphs_as_functors/index/dots}.

  The requirement that the one of the paths is nontrivial, however, is crucial in \cref{def:equalizers}.
\end{definition}

\begin{definition}\label{def:small_diagram}

\end{definition}

\begin{remark}\label{rem:inverting_isomorphisms_may_preserve_commutativity}
  Inverting isomorphisms in a \hyperref[def:categorical_diagram]{commutative diagram} may or may not preserve commutativity.

  If \( p = (f_1, \ldots, f_n) \) and \( q = (g_1, \ldots, g_m) \) are two paths in a commutative diagram, and if \( f_1 \) is invertible, then obviously
  \begin{equation*}
    f_n \bincirc \cdots \bincirc f_1 = g_m \bincirc \cdots \bincirc g_1
  \end{equation*}
  if and only if
  \begin{equation*}
    f_n \bincirc \cdots \bincirc f_2 = g_m \bincirc \cdots \bincirc g_1 \bincirc f_1
  \end{equation*}
  and similarly if \( f_n \) is invertible.

  On the other hand, consider \hyperref[def:ordinal]{ordinals} in \( \cat{Set} \). Denote by \( \iota \) the inclusion maps and by \( f: \omega^2 \to \omega \) the bijective map from \cref{thm:omega_equinumerous_with_omega_squared}. Then the following diagram commutes:
  \begin{equation}\label{eq:rem:inverting_isomorphisms_may_preserve_commutativity/ordinals_commuting}
    \begin{aligned}
      \includegraphics[page=1]{output/rem__inverting_isomorphisms_may_preserve_commutativity}
    \end{aligned}
  \end{equation}
  but the following does not:
  \begin{equation}\label{eq:rem:inverting_isomorphisms_may_preserve_commutativity/ordinals_not_commuting}
    \begin{aligned}
      \includegraphics[page=2]{output/rem__inverting_isomorphisms_may_preserve_commutativity}
    \end{aligned}
  \end{equation}
\end{remark}

\begin{definition}\label{def:functor_invertibility}
  In connection with \cref{def:morphism_invertibility} and \cref{def:function_invertibility}, we introduce the following terminology:
  \begin{thmenum}
    \thmitem{def:functor_invertibility/injective_on_objects} The \hyperref[def:functor]{functor} \( F: \cat{C} \to \cat{D} \) is \term{injective on objects} if the \hyperref[def:set_valued_map/restriction]{restriction}
    \begin{equation*}
      F\restr_{\obj(C)}: \obj(C) \to \obj(D)
    \end{equation*}
    is \hyperref[def:function_invertibility/injective]{injective}.

    That is, for every pair of objects \( A \) and \( B \) in \( \cat{C} \), from \( F(A) = F(B) \) it follows that \( A = B \).

    If, instead, from \( F(A) \cong F(B) \) it follows that \( A \cong B \), we say that \( F \) if \term{essentially injective on objects}.

    \thmitem{def:functor_invertibility/injective_on_morphisms} The \hyperref[def:functor]{functor} \( F: \cat{C} \to \cat{D} \) is \term{injective on morphisms} if its restriction to the set
    \begin{equation*}
      \bigcup\set{ \cat{C}(A, B) \given A, B \in \obj(\cat{C}) }
    \end{equation*}
    of all morphisms is injective.

    That is, for every pair of morphisms \( f \) and \( g \) in \( \cat{C} \), from \( F(f) = F(g) \) it follows that \( f = g \). Note that if the morphisms are not parallel, we assume that they are not equal.

    \thmitem{def:functor_invertibility/faithful}\mcite[def. 1.2.16]{Leinster2014BasicCategoryTheory} The functor \( F: \cat{C} \to \cat{D} \) is \term{faithful} if it is \hyperref[def:function_invertibility/injective]{injective} on \( \hom \)-sets, i.e. for all pairs of objects \( A \) and \( B \) in \( \cat{C} \), the restriction of \( F \) to \( \cat{C}(A, B) \) is an injective function.

    That is, for every pair of objects \( A \) and \( B \) in \( \cat{C} \) and every pair of morphisms \( f \) and \( g \) in \( \cat{C}(A, B) \), from \( F(f) = F(g) \) it follows that \( f = g \).

    See \cref{thm:def:functor_invertibility/injective} for how faithful functors relate to functors injective on objects or on morphisms.

    \thmitem{def:functor_invertibility/surjective_on_objects}\mcite[def. 1.3.17]{Leinster2014BasicCategoryTheory} The functor \( F: \cat{C} \to \cat{D} \) is \term{surjective on objects} if the restriction
    \begin{equation*}
      F\restr_{\obj(C)}: \obj(C) \to \obj(D)
    \end{equation*}
    is \hyperref[def:function_invertibility/surjective]{surjective}.

    That is, for every object \( B \) in \( \cat{D} \), there exists at least one object \( A \) in \( \cat{C} \) such that \( F(A) = B \).

    If, instead, there exists at least one object \( A \in \cat{C} \) such that \( F(A) \cong B \), we say that \( F \) is \term{essentially surjective on objects}.

    \thmitem{def:functor_invertibility/surjective_on_morphisms} Similarly, \( F: \cat{C} \to \cat{D} \) is \term{surjective on morphisms} if its restriction to the set of all morphisms is surjective.

    That is, for every morphism \( g \) in \( \cat{D} \), there exists at least one morphism \( f \) in \( \cat{C} \) such that \( F(f) = g \).

    \thmitem{def:functor_invertibility/full}\mcite[def. 1.2.16]{Leinster2014BasicCategoryTheory} The functor \( F: \cat{C} \to \cat{D} \) is \term{full} if it is surjective on \( \hom \)-sets, i.e. for all pairs of objects \( A \) and \( B \) in \( \cat{C} \), the restriction of \( F \) to \( \cat{C}(A, B) \) is a surjective function.

    That is, for every pair of objects \( A \) and \( B \) in \( \cat{C} \) and every morphism \( g: F(A) \to F(B) \) in \( \cat{D} \), there exists at least one morphism in \( f: A \to B \) in \( \cat{C} \) such that \( F(f) = g \).

    \thmitem{def:functor_invertibility/fully_faithful} Finally, \( F: \cat{C} \to \cat{D} \) is \term{fully faithful} if it is both full and faithful.
  \end{thmenum}
\end{definition}

\begin{proposition}\label{thm:commutative_diagrams_preserved_and_reflected}
  Functors preserve commutative diagrams and faithful functors also reflect commutative diagrams.

  More precisely, let \( \cat{C} \) be an arbitrary category, let \( D \) be a diagram in \( \cat{C} \), and let
  \begin{equation*}
    A \reloset {f_1} \to \anon \reloset {f_2} \to \cdots \reloset {f_{n-1}} \to \anon \reloset {f_n} \to B
  \end{equation*}
  and
  \begin{equation*}
    A \reloset {g_1} \to \anon \reloset {g_2} \to \cdots \reloset {g_{n-1}} \to \anon \reloset {g_m} \to B
  \end{equation*}
  be two \hyperref[def:graph_walk/directed]{directed walks} with the same endpoints in \( D \).

  For any functor \( F: \cat{C} \to \cat{D} \), if
  \begin{equation}\label{eq:thm:commutative_diagrams_preserved_and_reflected/source}
    f_n \bincirc \cdots \bincirc \bincirc f_1 = g_m \bincirc \cdots \bincirc \bincirc g_1,
  \end{equation}
  then
  \begin{equation}\label{eq:thm:commutative_diagrams_preserved_and_reflected/image}
    F(f_n) \bincirc \cdots \bincirc \bincirc F(f_1) = F(g_m) \bincirc \cdots \bincirc \bincirc F(g_1),
  \end{equation}

  Conversely, if \( F \) is faithful, then \eqref{eq:thm:commutative_diagrams_preserved_and_reflected/image} implies \eqref{eq:thm:commutative_diagrams_preserved_and_reflected/source}.
\end{proposition}
\begin{proof}
  Functors preserve composition by \ref{eq:def:functor/CF2}, hence \eqref{eq:thm:commutative_diagrams_preserved_and_reflected/image} follows from \eqref{eq:thm:commutative_diagrams_preserved_and_reflected/source} directly.

  Now suppose that \eqref{eq:thm:commutative_diagrams_preserved_and_reflected/image} holds for a faithful functor \( F \). \ref{eq:def:functor/CF2} allows us to reduce \eqref{eq:thm:commutative_diagrams_preserved_and_reflected/image} to
  \begin{equation*}
    F(f_n \bincirc \cdots \bincirc \bincirc f_1) = F(g_m \bincirc \cdots \bincirc \bincirc g_1).
  \end{equation*}

  Then, by injectivity of \( F \) on the \( \hom \)-set \( \cat{C}(\dom(f_1), \co\dom(f_1)) \), \eqref{eq:thm:commutative_diagrams_preserved_and_reflected/source} holds.
\end{proof}

\begin{proposition}\label{thm:def:functor_invertibility}
  \hyperref[def:functor]{Functors} have the following basic properties regarding their \hyperref[def:functor_invertibility]{invertibility}:

  \begin{thmenum}
    \thmitem{thm:def:functor_invertibility/injective} A functor is \hyperref[def:functor_invertibility/injective_on_morphisms]{injective on morphisms} if and only if it is both \hyperref[def:functor_invertibility/injective_on_objects]{injective on objects} and \hyperref[def:functor_invertibility/faithful]{faithful}.

    \thmitem{thm:def:functor_invertibility/surjective} A functor is \hyperref[def:functor_invertibility/surjective_on_morphisms]{surjective on morphisms} if and only if it is both \hyperref[def:functor_invertibility/surjective_on_objects]{surjective on objects} and \hyperref[def:functor_invertibility/full]{full}.

    \thmitem{thm:def:functor_invertibility/full_subcategory} A \hyperref[def:subcategory]{subcategory} \( \cat{D} \) of \( \cat{C} \) is full in the sense of \cref{def:subcategory} if and only if the \hyperref[def:subcategory]{inclusion functor} \( \Iota: \cat{D} \to \cat{C} \) is full in the sense of \cref{def:functor_invertibility/full}.

    \thmitem{thm:def:functor_invertibility/preserves_inverses} Any functor preserves \hyperref[def:morphism_invertibility/left]{left}, \hyperref[def:morphism_invertibility/right]{right inverses} and \hyperref[def:morphism_invertibility/isomorphism]{two-sided inverses}.

    \thmitem{thm:def:functor_invertibility/faithful_reflects_composition} \hyperref[def:functor_invertibility/faithful]{Faithful} functors reflect composition. That is, for every functor \( F: \cat{C} \to \cat{D} \), if the following diagram commutes:
    \begin{equation}\label{eq:thm:def:functor_invertibility/faithful_reflects_composition/image}
      \begin{aligned}
        \includegraphics[page=1]{output/thm__def__functor_invertibility__properties}
      \end{aligned}
    \end{equation}
    then the following diagram are identities:
    \begin{equation}\label{eq:thm:def:functor_invertibility/faithful_reflects_composition/source}
      \begin{aligned}
        \includegraphics[page=2]{output/thm__def__functor_invertibility__properties}
      \end{aligned}
    \end{equation}

    \thmitem{thm:def:functor_invertibility/faithful_reflects_cancellative} A \hyperref[def:functor_invertibility/faithful]{faithful} functor reflects monomorphisms and epimorphisms. That is, for every functor \( F: \cat{C} \to \cat{D} \) and morphism \( f: A \to B \) in \( \cat{C} \), if \( F(f) \) is a monomorphism (resp. epimorphism), so is \( f \).

    \thmitem{thm:def:functor_invertibility/fully_faithful_reflects_identities} A \hyperref[def:functor_invertibility/fully_faithful]{fully faithful} functor reflects identities. That is, for every functor \( F: \cat{C} \to \cat{D} \) and endomorphism \( f: A \to A \) in \( \cat{C} \), if \( F(f) = \id_{F(A)} \), then \( f = \id_A \).

    \thmitem{thm:def:functor_invertibility/fully_faithful_reflects_invertible} A \hyperref[def:functor_invertibility/fully_faithful]{fully faithful} functor reflects split monomorphisms and split epimorphisms.

    That is, for every functor \( F: \cat{C} \to \cat{D} \) and morphism \( f: A \to B \) in \( \cat{C} \), if \( F(f) \) is a split monomorphism (resp. split epimorphism or isomorphism), so is \( f \).

    \thmitem{thm:def:functor_invertibility/isomorphism} A functor between \hyperref[def:small_set]{small} categories that is both injective and surjective on morphisms is itself an isomorphism in \( \cat{Cat} \).
  \end{thmenum}
\end{proposition}
\begin{proof}
  \SubProofOf{thm:def:functor_invertibility/injective}
  \SufficiencySubProof* Let \( F: \cat{C} \to \cat{D} \) be injective on morphisms. It is trivially faithful since faithfulness is a more restrictive condition.

  To see that \( F \) is injective on objects, let \( A, B \in \cat{C} \) and suppose that \( F(A) = F(B) \). Then \( \id_{F(A)} = \id_{F(B)} \) and
  \begin{equation*}
    F(\id_A)
    \reloset {\ref{def:functor/CF2}} =
    \id_{F(A)}
    =
    \id_{F(B)}
    \reloset {\ref{def:functor/CF2}} =
    F(\id_B).
  \end{equation*}

  Since \( F \) is injective on morphisms, it follows that \( \id_A = \id_B \), hence \( A = B \). Thus, \( F \) is injective on objects.

  \NecessitySubProof* Let \( F: \cat{C} \to \cat{D} \) be faithful and injective on objects. Let \( f: A \to B \) and \( g: C \to D \) be morphisms in \( \cat{C} \) such that \( F(f) = F(g) \).

  Then both \( F(f) \) and \( F(g) \) have the same domain \( F(A) = F(C) \) and codomain \( F(B) = F(D) \). Hence, since \( F \) is injective on objects, we have \( A = C \) and \( B = D \).

  Thus, \( f \) and \( g \) are both morphisms from \( A \) to \( B \). Since \( F \) is also faithful, from \( F(f) = F(g) \) it follows that \( f = g \).

  Therefore, \( F \) is injective on morphisms.

  \SubProofOf{thm:def:functor_invertibility/surjective}
  \SufficiencySubProof* Let \( F: \cat{C} \to \cat{D} \) be surjective on morphisms. It is trivially full since fullness is a more restrictive condition.

  To see that \( F \) is surjective on objects, let \( C \in \cat{D} \). Then there exists some morphism \( f: A \to B \) in \( \cat{C} \) such that \( F(f) = \id_Z \). We thus necessarily have \( F(A) = C \) and \( F(B) = C \).

  \NecessitySubProof* Let \( F: \cat{C} \to \cat{D} \) be full and injective on objects. Let \( g: C \to D \) be a morphism in \( \cat{D} \).

  Since \( F \) is surjective on objects, there exists preimages \( A \) of \( C \) and \( B \) of \( D \) under \( F \). Thus, \( g \in \cat{D}(F(A), F(B)) \).

  Since \( F \) is also full, there exists some morphism \( f: A \to B \) such that \( F(f) = g \).

  Therefore, \( F \) is surjective on morphisms.

  \SubProofOf{thm:def:functor_invertibility/full_subcategory} Trivial.

  \SubProofOf{thm:def:functor_invertibility/preserves_inverses} Let \( g: B \to A \) be a left inverse of \( f: A \to B \). Then
  \begin{equation*}
    F(g) \bincirc F(f)
    \reloset {\ref{def:functor/CF2}} =
    F(g \bincirc f)
    =
    F(\id_A)
    \reloset {\ref{def:functor/CF1}} =
    \id_{F(A)}.
  \end{equation*}

  The case of right inverses is similar.

  \SubProofOf{thm:def:functor_invertibility/faithful_reflects_composition} Suppose that \eqref{eq:thm:def:functor_invertibility/faithful_reflects_composition/image} commutes. Then, since \( F \) is faithful and thus injective on the \( \hom \)-set \( \cat{C}(A, C) \), the equality \( F(g \bincirc f) = F(g) \bincirc F(f) = F(h) \) implies that \( g \bincirc f = h \). Hence, \eqref{eq:thm:def:functor_invertibility/faithful_reflects_composition/source} also commutes.

  \SubProofOf{thm:def:functor_invertibility/faithful_reflects_cancellative} Let \( F(g) \) be a monomorphism and let \( f_1, f_2: A \to B \) be parallel morphisms such that
  \begin{equation*}
    g \bincirc f_1 = g \bincirc f_2.
  \end{equation*}

  Then, since \( F(g) \) is a monomorphism, we have that \( F(f_1) = F(f_2) \). Since \( F \) is faithful, the restriction \( F\restr_{C(A, B)} \) is injective, and \( f_1 = f_2 \).

  The proof when \( F(g) \) is an epimorphism is analogous.

  \SubProofOf{thm:def:functor_invertibility/fully_faithful_reflects_identities} If \( F: \cat{C} \to \cat{D} \) is fully faithful, for every object \( A \) in \( \cat{C} \), the identity morphism \( \id_{F(A)} \) has a unique preimage under \( F \). By \ref{def:functor/CF1}, this preimage can only be \( \id_A \).

  \SubProofOf{thm:def:functor_invertibility/fully_faithful_reflects_invertible} Let \( q \) be a left inverse of \( F(f) \). Since \( F \) is fully faithful, there exists a unique morphism \( g: B \to A \) such that \( F(g) = q \).

  Since
  \begin{equation*}
    F(g) \bincirc F(f) = \id_{F(A)},
  \end{equation*}
  by \cref{thm:def:functor_invertibility/fully_faithful_reflects_identities} we have
  \begin{equation*}
    g \bincirc f = \id_A.
  \end{equation*}

  Therefore, \( g \) is a left inverse of \( F(f) \).

  The proof for right inverses follows from \cref{thm:def:morphism_invertibility/inverse_duality}.

  From \cref{thm:def:morphism_invertibility/left_and_right} it follows that if \( F(f) \) is an isomorphism, so is \( f \).

  \SubProofOf{thm:def:functor_invertibility/isomorphism} If \( F \) is both injective and surjective on morphisms, it is also injective and surjective on objects and hence, as a function, is bijective. Therefore, it is both left and right invertible as a consequence of \cref{thm:function_invertibility_categorical/fully_invertible}.
\end{proof}

\begin{example}\label{ex:def:functor_invertibility}
  \hfill
  \begin{thmenum}
    \thmitem{ex:def:functor_invertibility/power} The power set functor described in \cref{ex:def:functor} is clearly \hyperref[def:functor_invertibility/injective_on_morphisms]{injective on morphisms}, hence by \cref{thm:def:functor_invertibility/injective}, it is also \hyperref[def:functor_invertibility/injective_on_objects]{injective on objects} and \hyperref[def:functor_invertibility/faithful]{faithful}.

    It is not full, nor surjective on objects.

    \thmitem{ex:def:functor_invertibility/cat_to_set} The forgetful functor \( D: \cat{Cat} \to \cat{Set} \) discussed in \cref{def:discrete_category} is \hyperref[def:functor_invertibility/surjective_on_morphisms]{surjective on morphisms}, hence by \cref{thm:def:functor_invertibility/surjective}, it is also \hyperref[def:functor_invertibility/surjective_on_objects]{surjective on objects} and \hyperref[def:functor_invertibility/full]{full}.

    It is not faithful, nor injective on objects.
  \end{thmenum}
\end{example}

\begin{definition}\label{def:natural_transformation}\mcite[def. 1.3.1]{Leinster2014BasicCategoryTheory}
  Let \( F \) and \( G \) be parallel \hyperref[def:functor]{functors} from the category \( \cat{C} \) to \( \cat{D} \).

  A \term{natural transformation} \( \alpha \) from \( F \) to \( G \) is an \hyperref[def:indexed_family]{indexed family} of
  \begin{equation}\label{eq:def:natural_transformation/family}
    \seq{ \alpha_A: F(A) \to G(A) }_{A \in \cat{C}}
  \end{equation}
  of morphisms in \( \cat{D} \) such that, for every morphism \( f: A \to B \) in \( \cat{C} \), the following \hyperref[def:categorical_diagram]{diagram commutes}:
  \begin{equation}\label{eq:def:natural_transformation/diagram}
    \begin{aligned}
      \includegraphics[page=1]{output/def__natural_transformation}
    \end{aligned}
  \end{equation}

  The morphisms \( \alpha_A \) are called the components of \( \alpha \). We denote natural transformations by \( \alpha: F \Rightarrow G \) and, when used in diagrams, by
  \begin{equation}\label{eq:def:natural_transformation/notation}
    \begin{aligned}
      \includegraphics[page=2]{output/def__natural_transformation}
    \end{aligned}
  \end{equation}
\end{definition}

\begin{example}\label{ex:directed_multigraphs_as_functors}\mcite[example 1.3.46]{Perrone2024StartingCategoryTheory}
  In \cref{def:directed_multigraph}, we have defined a directed multigraph as a set \( V \) of vertices, a set \( A \) of arcs and two functions --- the initial \( i: A \to V \) and terminal \( t: A \to V \) vertices of an arc.

  Now consider the following \hyperref[def:categorical_diagram]{index category} \( \cat{I}: \)
  \begin{equation}\label{eq:ex:directed_multigraphs_as_functors/index/dots}
    \begin{aligned}
      \includegraphics[page=1]{output/ex__directed_multigraphs_as_functors}
    \end{aligned}
  \end{equation}

  For the sake of readability, we will give the following explicit labels in this category:
  \begin{equation}\label{eq:ex:directed_multigraphs_as_functors/index/annotated}
    \begin{aligned}
      \includegraphics[page=2]{output/ex__directed_multigraphs_as_functors}
    \end{aligned}
  \end{equation}

  A directed multigraph can then be defined as a functor \( G: \cat{I} \to \cat{Set} \) to the category \hyperref[def:category_of_small_sets]{\( \cat{Set} \)} of \hyperref[def:small_set]{small} sets.

  A \hyperref[def:natural_transformation]{natural transformation} from the directed multigraph \( G: \cat{I} \to \cat{Set} \) to \( H: \cat{I} \to \cat{Set} \) is then a pair of functions \( f_V: G(V) \to H(V) \) and \( f_A: G(A) \to H(A) \) such that the following \hyperref[def:categorical_diagram]{diagrams commute}:
  \begin{equation}\label{eq:ex:directed_multigraphs_as_functors/index/diagram}
    \begin{aligned}
      \includegraphics[page=3]{output/ex__directed_multigraphs_as_functors}
      \quad\quad\quad\quad
      \includegraphics[page=4]{output/ex__directed_multigraphs_as_functors}
    \end{aligned}
  \end{equation}

  See \cref{ex:isomorphism_of_directed_multigraph_categories} for how these functors related to directed multigraphs as defined in \cref{def:directed_multigraph}.
\end{example}

\begin{remark}\label{rem:natural_transformations_into_set}
  Let \( \cat{C} \) be an arbitrary \hyperref[def:small_set]{small} category. A \hyperref[def:natural_transformation]{natural transformation} \( \alpha \) from \( F: \cat{C} \to \cat{Set} \) to \( G: \cat{C} \to \cat{Set} \) is then a family of functions
  \begin{equation*}
    \seq{ \alpha_A: F(A) \to G(A) }_{A \in \cat{C}}.
  \end{equation*}

  Suppose that for every two objects \( A \) and \( B \) in \( \cat{C} \), the functions \( \alpha_A \) and \( \alpha_B \) agree on \( F(A) \cap F(B) \). This is automatically satisfied in \( F(A) \) and \( F(B) \) are disjoint whenever \( A \neq B \).

  We can then take the set-theoretic union of \( \alpha \) to obtain the function
  \begin{equation*}
    \bigcup_{A \in \cat{C}} \alpha_A: \bigcup\set{ F(A) \given A \in \cat{C} } \to \bigcup\set{ G(A) \given A \in \cat{C} }.
  \end{equation*}

  Both the domain and codomain are sets as a consequence of \ref{def:grothendieck_universe/union}, therefore the function is well-defined in the \hyperref[def:smallest_grothendieck_universe]{smallest Grothendieck universe} \( \mscrU_0 \). Denote it on \( \Alpha \) for brevity.

  An advantage of this is that we can define a natural transformation to be a function on a general enough set and then prove that its restrictions satisfy \eqref{eq:def:natural_transformation/diagram}.

  For example, consider the power set functor \( \pow: \cat{Set} \to \cat{Set} \) discussed in \cref{ex:def:functor}. The \hyperref[def:set_valued_map/identity]{identity function} \( \id_{\mscrU_0} \) is then a natural transformation from the identity functor \( \id_{\cat{Set}} \) to \( \pow \).

  Another natural transformation between the same functors is the singleton set operation \( \Sigma \) on sets defined as \( A \mapsto \set{ A } \). Note that, in this context, \( \Sigma \) operates not on the sets \( \id_{\cat{Set}}(A) \) and \( \pow(A) \), but on their members. The diagram \eqref{eq:def:natural_transformation/diagram} becomes
  \begin{equation}\label{eq:rem:natural_transformations_into_set}
    \begin{aligned}
      \includegraphics[page=1]{output/rem__natural_transformations_into_set}
    \end{aligned}
  \end{equation}

  This diagram commutes because, for every function \( f: A \to B \) and every \( x \in A \), we have
  \begin{equation*}
    f[\set{ x }] = \set{ f(x) }.
  \end{equation*}
\end{remark}

\begin{definition}\label{def:functor_category}
  Let \( \cat{C} \) and \( \cat{D} \) be arbitrary \hyperref[def:category]{categories}. The \term{functor category} \( [\cat{C}, \cat{D}] \), also denoted as \( \cat{D}^{\cat{C}} \), is defined as follows:

  \begin{itemize}
    \item The \hyperref[def:category/objects]{set of objects} \( \obj([\cat{C}, \cat{D}]) \) is the set of all functors from \( \cat{C} \) to \( \cat{D} \).

    \item The \hyperref[def:category/morphisms]{set of morphisms} \( [\cat{C}, \cat{D}](F, G) \) from \( F \) to \( G \) is the set of all \hyperref[def:natural_transformation]{natural transformations} from \( F \) to \( G \).

    \item The \hyperref[def:category/composition]{composition of the morphisms} \( \alpha: F \Rightarrow G \) and \( \beta: G \Rightarrow H \) is the natural transformation \( \beta \bincirc \alpha: F \Rightarrow H \) defined in terms of componentwise morphism composition, i.e.
    \begin{equation}\label{eq:def:functor_category/composition}
      (\beta \bincirc \alpha)_A \coloneqq \beta_A \bincirc \alpha_A.
    \end{equation}

    \item The \hyperref[def:category/identity]{identity morphism} on the functor \( F: \cat{C} \to \cat{D} \) is the \term{identity natural transformation} \( \id_F: F \Rightarrow F \) with components
    \begin{equation}\label{eq:def:functor_category/identity}
      (\id_F)_A \coloneqq \underbrace{\id_{F(A)}}_{F(\id_A)}
    \end{equation}
  \end{itemize}
\end{definition}
\begin{defproof}
  Just to verify that the composition \( \beta \bincirc \alpha \) defined in \eqref{eq:def:functor_category/composition} is indeed a natural transformation from \( F \) to \( H \), note that the following diagram trivially commutes:
  \begin{equation}\label{def:functor_category/composition}
    \begin{aligned}
      \includegraphics[page=1]{output/def__functor_category}
    \end{aligned}
  \end{equation}

  Now, to see that \( [\cat{C}, \cat{D}] \) is indeed a category, we verify the conditions \ref{def:category/C1} and \ref{def:category/C2}, which are in turn inherited from the same conditions on the categories \( \cat{C} \) and \( \cat{D} \).

  \SubProofOf{def:category/C1} For every two functors \( F, G: \cat{C} \to \cat{D} \) and natural transformation \( \alpha: F \Rightarrow G \), for every object \( A \in \cat{C} \) we have
  \begin{equation*}
    \id_{G(A)} \bincirc \alpha_A
    \reloset{\eqref{def:category/C1}} =
    \alpha_A
    \reloset{\eqref{def:category/C1}} =
    \alpha_A \bincirc \id_{F(A)}
  \end{equation*}

  Therefore,
  \begin{equation*}
    \id_G \bincirc \alpha = \alpha = \alpha \bincirc \id_F
  \end{equation*}
  and, after generalizing, we obtain that \ref{def:category/C1} holds in \( [\cat{C}, \cat{D}] \).

  \SubProofOf{def:category/C2} For any quadruple \( F \), \( G \), \( H \) and \( T \) of functors from \( \cat{C} \) to \( \cat{D} \) and every combination of natural transformations \( \alpha: F \Rightarrow G \), \( \beta: G \Rightarrow H \) and \( \gamma: H \Rightarrow T \), for every object \( A \in \cat{C} \) we have
  \begin{equation*}
    (\gamma_A \bincirc \beta_A) \bincirc \alpha_A
    \reloset{\eqref{def:category/C2}} =
    \gamma_A \bincirc (\beta_A \bincirc \alpha_A).
  \end{equation*}

  Therefore, after generalizing, we obtain that \ref{def:category/C2} holds in \( [\cat{C}, \cat{D}] \).
\end{defproof}

\begin{example}\label{ex:isomorphism_of_directed_multigraph_categories}
  In \cref{ex:directed_multigraphs_as_functors}, we defined \hyperref[def:directed_multigraph]{directed multigraphs} as functors from a certain index category \( \cat{I} \) to \( \cat{Set} \) (for a \hyperref[def:small_category]{fixed Grothendieck universe} \( \mscrU \)).

  There is then an obvious correspondence between directed multigraphs as objects of the category \( \cat{DM} \) defined in \cref{def:directed_multigraph/category}, and directed multigraphs as objects in the \hyperref[def:functor_category]{functor category} \( [\cat{I}, \cat{Set}] \), defined in \cref{ex:directed_multigraphs_as_functors}. Indeed, given any functor \( G: \cat{I} \to \cat{Set} \), the quadruple
  \begin{equation*}
    \parens[\big]{ G(V), G(A), G(h), G(t) }
  \end{equation*}
  is a directed multigraph in the sense of \cref{def:directed_multigraph}.

  No object in \( \cat{DM} \) is formally equal to any object in \( [\cat{I}, \cat{Set}] \) in the sense of \hyperref[def:zfc]{\logic{ZFC}}. They are, however, equivalent, as shown above, and this can be formalized by stating that the two categories are isomorphic, in the sense of \cref{def:morphism_invertibility/isomorphism}, as objects of the category \( \cat[\mscrV]{Cat} \), where \( \mscrV \) is a Grothendieck universe that strictly contains \( \mscrU \). We have already defined this isomorphism explicitly.

  This is an example of \term{isomorphism of categories}. In practice, if two categories are not so obviously identical, we are usually better served by \term{equivalences of categories} defined in \cref{def:category_equivalence}.
\end{example}

\begin{definition}\label{def:opposite_natural_transformation}\mimprovised
  The \term{opposite} natural transformation of \( \alpha: F \Rightarrow G \), where \( F \) and \( G \) are functors from \( \cat{C} \) to \( \cat{D} \), is the natural transformation \( \alpha^\oppos: G^\oppos \Rightarrow F^\oppos \), in which we take the opposite of each component in \( \alpha \).

  Dual natural transformation arise naturally in \cref{thm:opposite_of_functor_category}.
\end{definition}
\begin{defproof}
  The naturality diagram \eqref{eq:def:natural_transformation/diagram} commutes for \( \alpha^\oppos \) because all morphisms are simply reversed.
\end{defproof}

\begin{definition}\label{def:opposite_functor}\mcite[33]{MacLane1998CategoryTheory}
  The \term{opposite} functor of \( F: \cat{C} \to \cat{D} \) is the functor
  \begin{equation*}
    \begin{aligned}
      &F^\oppos: \cat{C}^\oppos \to \cat{D}^\oppos \\
      &F^\oppos(A) \coloneqq A \\
      &F^\oppos(f^\oppos: B \to A) \coloneqq [F(f: A \to B)]^\oppos.
    \end{aligned}
  \end{equation*}

  For the composition of functors, we then have
  \begin{equation}\label{eq:def:opposite_functor/composition}
    [G \bincirc F]^\oppos = G^\oppos \bincirc F^\oppos.
  \end{equation}

  This is somewhat in contrast to the general practice of inverting morphisms when taking opposites. Thus, we have the \hyperref[def:contravariant_functor]{contravariant} \term{oppositization functor}
  \begin{equation*}
    (\anon*)^\oppos: \cat{Cat}^\oppos \to \cat{Cat}.
  \end{equation*}

  As an \hyperref[def:function/endofunction]{endofunction} on \( \obj(\cat{Cat}) \), the oppositization functor is clearly an \hyperref[def:categorical_involution]{involution}.

  Dual functors also arise naturally in \cref{thm:opposite_of_functor_category}.
\end{definition}

\begin{proposition}\label{thm:opposite_of_functor_category}
  For the \hyperref[def:opposite_category]{opposite} of the \hyperref[def:functor_category]{functor category} \( [\cat{C}, \cat{D}] \) we have
  \begin{equation*}
    [\cat{C}, \cat{D}]^\oppos = [\cat{C}^\oppos, \cat{D}^\oppos].
  \end{equation*}

  This is part of the duality principles listed in \fullref{thm:category_duality}.
\end{proposition}
\begin{proof}
  In \cref{def:opposite_functor}, we have defined the opposite functor \( F^\oppos: \cat{C}^\oppos \to \cat{D}^\oppos \) of \( F: \cat{C} \to \cat{D} \) in a way that allows us to regard it as an object of \( [\cat{C}^\oppos, \cat{D}^\oppos] \).

  In \cref{def:opposite_natural_transformation}, we have defined the opposite natural transformation \( \alpha^\oppos: G^\oppos \to F^\oppos \) of \( \alpha: F \Rightarrow G \) in a way that allows us to regard it as a morphism of \( [\cat{C}^\oppos, \cat{D}^\oppos] \).

  Furthermore, \( \alpha^\oppos: G^\oppos \to F^\oppos \) reverses the direction of its morphisms, and hence it is the dual to \( \alpha \) in the category \( [\cat{C}, \cat{D}]^\oppos \).
\end{proof}

\begin{definition}\label{def:diagonal_functor}\mcite[exerc. 3.1.1]{Leinster2014BasicCategoryTheory}
  Given an \term{index category} \( \cat{I} \) and an arbitrary category \( \cat{C} \), for any object \( A \) in \( \cat{C} \), we can define the \term{constant functor}
  \begin{equation*}
    \begin{aligned}
      &\Delta_A^{\cat{I}}: \cat{I} \to \cat{C}, \\
      &\Delta_A^{\cat{I}}(X) \coloneqq X, \\
      &\Delta_A^{\cat{I}}(g: X \to Y) \coloneqq \id_A.
    \end{aligned}
  \end{equation*}

  Given two objects \( A \) and \( B \) in \( \cat{C} \), a natural transformation \( \alpha: \Delta_A^{\cat{I}} \Rightarrow \Delta_B^{\cat{I}} \) is an \hyperref[def:indexed_family]{indexed family} that gives the same morphism for every object of the index category \( \cat{I} \).

  Indeed, the diagram \eqref{eq:def:natural_transformation/diagram} in this case becomes
  \begin{equation}\label{eq:def:diagonal_functor/nat}
    \begin{aligned}
      \includegraphics[page=1]{output/def__diagonal_functor}
    \end{aligned}
  \end{equation}

  This diagram implies that \( \alpha_A = \alpha_B \) for any two objects \( A \) and \( B \) in \( \cat{I} \). Therefore, all components of \( \alpha \) are equal to some morphism in \( \cat{C}(A, B) \).

  We can now define the \( \cat{I} \)-shaped \term{diagonal functor} on \( \cat{C} \)
  \begin{equation*}
    \begin{aligned}
      &\Delta^{\cat{I}}: \cat{C} \to [\cat{I}, \cat{C}], \\
      &\Delta^{\cat{I}}(A) \coloneqq \Delta_A^{\cat{I}}, \\
      &\Delta^{\cat{I}}(f: A \to B) \coloneqq \seq{ f: A \to B }_{k \in \cat{I}}.
    \end{aligned}
  \end{equation*}

  It is called a diagonal functor because, if \( \cat{I} \) is a discrete category of small two objects, then \( \Delta_{\cat{I}} \) gives the diagonal of the \hyperref[def:product_category]{product category} \( \cat{C}^2 \) by providing, for each object \( A \) of \( \cat{C} \), the ordered pair \( (A, A) \) (and similarly for morphisms).
\end{definition}

\begin{proposition}\label{thm:natural_isomorphism}\mcite{math3ma:natural_transformations}
  Let \( F \) and \( G \) be parallel \hyperref[def:functor]{functors} from the category \( \cat{C} \) to \( \cat{D} \). The family \eqref{eq:def:natural_transformation/family} is an isomorphism in the corresponding \hyperref[def:functor_category]{functor category} \( [\cat{C}, \cat{D}] \) if and only if all of its components are isomorphisms and, for any morphism \( f: A \to B \) in \( \cat{C} \), the following diagram commutes:
  \begin{equation}\label{eq:thm:natural_isomorphism/diagram}
    \begin{aligned}
      \includegraphics[page=1]{output/thm__natural_isomorphism}
    \end{aligned}
  \end{equation}

  We say that \( \alpha \) is a \term{natural isomorphism}.
\end{proposition}
\begin{proof}
  If all components of \( \alpha \) are isomorphisms, the condition
  \begin{equation*}
    \alpha_B \bincirc F(f) = G(f) \bincirc \alpha_A
  \end{equation*}
  is equivalent to
  \begin{equation*}
    F(f) = \alpha_B^{-1} \bincirc G(f) \bincirc \alpha_A.
  \end{equation*}

  We must now show that, if \( \alpha \) is an isomorphism in \( [\cat{C}, \cat{D}] \), all of its components are isomorphisms.

  If \( \alpha: F \Rightarrow G \) is an isomorphism in \( [\cat{C}, \cat{D}] \). Then there exists some natural transformation \( \beta: G \Rightarrow F \) such that
  \begin{equation*}
    \beta \bincirc \alpha = \id_F \quad\T{and}\quad \alpha \bincirc \beta = \id_G.
  \end{equation*}

  For every object \( A \) in \( \cat{C} \), the morphism \( \alpha_A: F(A) \to G(A) \) composed with \( \beta_A: G(A) \to F(A) \) is
  \begin{equation*}
    \beta_A \bincirc \alpha_A = \id_{F(A)}.
  \end{equation*}

  Therefore, \( \alpha_A \) is left invertible. Analogously,
  \begin{equation*}
    \alpha_A \bincirc \beta_A = \id_{F(A)}
  \end{equation*}
  and hence \( \alpha_A \) is right invertible.

  Therefore, for every object \( A \) in \( \cat{C} \), the morphism \( \alpha_A \) is fully invertible, i.e. an isomorphism.
\end{proof}

\begin{definition}\label{def:product_category}\mcite[const. 1.1.11]{Leinster2014BasicCategoryTheory}
  We define the \term{product category} \( \cat{C} \times \cat{D} \) of \( \cat{C} \) and \( \cat{D} \) as follows:

  \begin{itemize}
    \item The \hyperref[def:category/objects]{set of objects} is the \hyperref[def:cartesian_product]{Cartesian product}
    \begin{equation}\label{eq:def:product_category/objects}
      \obj(\cat{C} \times \cat{D}) \coloneqq \obj(\cat{C}) \times \obj(\cat{D}).
    \end{equation}

    \item The \hyperref[def:category/morphisms]{set of morphisms} from the pair of objects \( (A, X) \) to \( (B, Y) \) is the product
    \begin{equation}\label{eq:def:product_category/morphisms}
      (\cat{C} \times \cat{D})\parens[\big]{ (A, X), (B, Y) } \coloneqq \cat{C}(A, B) \times \cat{D}(X, Y).
    \end{equation}

    \item The \hyperref[def:category/composition]{composition of the morphisms}
    \begin{align*}
      (f, r)&: (A, X) \to (B, Y) \\
      (g, s)&: (B, Y) \to (C, Z)
    \end{align*}
    is the pairwise composition
    \begin{equation}\label{eq:def:product_category/composition}
      (g, s) \bincirc (f, r) \coloneqq \underbrace{(g \bincirc f, s \bincirc r)}_{(A, X) \to (C, Z)}.
    \end{equation}

    \item The \hyperref[def:category/identity]{identity morphism} of the pair \( (A, X) \) is simply the pair of identity morphisms \( (\id_A, \id_X) \).
  \end{itemize}
\end{definition}

\begin{definition}\label{def:hom_functor}
  Let \( \cat{C} \) be a \hyperref[def:small_category]{locally \hyperref[def:small_set]{small}} category. We can regard the \( \hom \)-sets \( \cat{C}(A, B) \) as a functor parameterized by objects of \( \cat{C} \).

  \begin{thmenum}
    \thmitem{def:hom_functor/binary} For any pair of morphisms \( f: B \to A \) and \( g: X \to Y \) in \( \cat{C} \), define the operator
    \begin{equation}\label{eq:def:hom_functor/t}
      \begin{aligned}
        &T_{f,g}: \cat{C}(A, X) \to \cat{C}(B, Y) \\
        &T_{f,g}(s) \mapsto g \bincirc s \bincirc f.
      \end{aligned}
    \end{equation}

    The action of \( T_{f,g} \) can be expressed graphically as
    \begin{equation}\label{eq:def:hom_functor/t_diagram}
      \begin{aligned}
        \includegraphics[page=1]{output/def__hom_functor}
      \end{aligned}
    \end{equation}

    We can now define the following \term{binary hom-functor}:
    \begin{equation}\label{eq:def:hom_functor/binary}
      \begin{aligned}
        &\cat{C}(\anon*, \anon*): \cat{C}^\oppos \times \cat{C} \to \cat{Set} \\
        &\cat{C}(A, X) \coloneqq \set{ s: A \to X } \\
        &\cat{C}(f, g) \coloneqq T_{f,g}
      \end{aligned}
    \end{equation}

    \thmitem{def:hom_functor/unary} Fixing the first argument \( A \) in \eqref{eq:def:hom_functor/binary}, we instead obtain a covariant unary hom-functor:
    \begin{equation}\label{eq:def:hom_functor/unary/covariant}
      \cat{C}(A, \anon*): \cat{C} \to \cat{Set}
    \end{equation}

    Analogously, fixing the second argument \( X \), we obtain a \hyperref[def:hom_functor/unary]{contravariant} unary hom-functor:
    \begin{equation}\label{eq:def:hom_functor/unary/contravariant}
      \cat{C}(\anon*, X): \cat{C}^\oppos \to \cat{Set}
    \end{equation}
  \end{thmenum}
\end{definition}
\begin{defproof}
  It is sufficient to verify that \eqref{eq:def:hom_functor/binary} defines a functor. \ref{def:functor/CF2} can be seen to hold by inspecting the diagram:
  \begin{equation}\label{eq:def:hom_functor/inv_composition}
    \begin{aligned}
      \includegraphics[page=2]{output/def__hom_functor}
    \end{aligned}
  \end{equation}

  The other functor condition \ref{def:functor/CF1} is straightforward to prove.
\end{defproof}

\begin{definition}\label{def:function_currying}\mcite[259]{TroelstraSchwichtenberg2000BasicProofTheory}
  Given the \hyperref[con:function_arguments]{binary} function \( f: A \times B \to C \), we can define another function \( g: A \to \fun(B, C) \) as
  \begin{equation*}
    g(x) \coloneqq (y \mapsto f(x, y)).
  \end{equation*}

  We call this process \term{currying} after Haskell Curry. Obtaining \( f \) from \( g \) can be called \term{uncurrying}.
\end{definition}
\begin{comments}
  \item Currying is useful if we have somehow fixed a value \( x_0 \in A \), in which case we can \enquote{get rid} of one argument by introducing some shortcut for the function \( g(x_0): B \to C \) for the sake of reducing notational clutter. See \cref{def:differentiability/first_variation} and our proof of \cref{thm:countably_infinite_union_of_countably_infinite_sets} for example of how this is useful in the wild.
\end{comments}

\begin{proposition}\label{thm:currying_is_natural_isomorphism}
  \hyperref[def:function_currying]{Function currying} is a natural isomorphism between the functors
  \begin{align*}
    &\cat{Set}(A \times B, C)
    &\cat{Set}(A, \cat{Set}(B, C))
  \end{align*}

  More concretely, consider the following functors, which are loosely based on \cref{def:hom_functor}:
  \begin{equation*}
    \begin{aligned}
      &V: \cat{Set}^3 \to \cat{Set} \\
      &V(A, B, C) \coloneqq \cat{Set}(A \times B, C) \\
      \Big[& V(f: X \to A, g: Y \to B, h: C \to Z) \Big](s: A \times B \to C) \coloneqq \smash{ \overbrace{ (x, y) \mapsto h\parens[\Bigg]{ s\parens[\big]{ \underbrace{ f(x), g(y) }_{A \times B} } } }^{X \times Y \to Z} } \\
    \end{aligned}
  \end{equation*}
  and
  \begin{equation*}
    \begin{aligned}
      &W: \cat{Set}^3 \to \cat{Set} \\
      &W(A, B, C) \coloneqq \cat{Set}(A, \cat{Set}(B, C)) \\
      \Big[& W(f: X \to A, g: Y \to B, h: C \to Z) \Big](t: A \to \cat{Set}(B, C)) \coloneqq \smash{ \underbrace{ x \mapsto \overbrace{y \mapsto h\parens[\Bigg]{ \overbrace{t(f(x))}^{B \to C}\parens[\big]{ g(y) } } }^{Y \to Z} }_{X \to \cat{Set}(Y, Z)} } \\
    \end{aligned}
  \end{equation*}

  Then the family of functions
  \begin{equation*}
    \begin{aligned}
      &\alpha: V \Rightarrow W \\
      &\alpha_{A,B,C}(s: A \times B \to C) \coloneqq a \mapsto b \mapsto s(a, b)
    \end{aligned}
  \end{equation*}
  is a \hyperref[thm:natural_isomorphism]{natural isomorphism}.
\end{proposition}
\begin{proof}
  The function \( \varphi \) is clearly invertible. Fix a triple of functions \( f: X \to A \), \( g: Y \to B \) and \( h: C \to Z \). For every \( s: A \times B \to C \) we have
  \begin{align*}
    [W(f, g, h)](\alpha_{A,B,C}(s))
    &=
    x \mapsto y \mapsto \parens[\big]{ h\parens[\big]{ [a \mapsto b \mapsto s(a, b)](f(x))(g(y)) } }
    = \\ &=
    x \mapsto y \mapsto h\parens[\big]{ s(f(x), g(y)) }
    = \\ &=
    \alpha_{X,Y,Z}(V(f, g, h)),
  \end{align*}
  which proves that the following diagram commutes:
  \begin{equation}\label{eq:thm:currying_is_natural_isomorphism/diagram}
    \begin{aligned}
      \includegraphics[page=1]{output/thm__currying_is_natural_isomorphism}
    \end{aligned}
  \end{equation}
\end{proof}
